\section{Conclusion}

\textbf{What was built.} We built and audited an end-to-end, AI-guided CPT investigation loop. It combines
five AI predictors, an uncertainty--disagreement--distance acquisition score, a multi-signal stop rule and
an owner's economic model.

\textbf{Findings.} (i)~the TabPFN hybrids gave the lowest profile error (0.497 log-MAE) and the highest low-cell recall (36\%), within about 3\% of XGBoost and GP; (ii)~with rules frozen on development sites, no early stop failed the audit; (iii)~CPTs were saved only at the most uniform site: 47.6\%, at the minimum the rule allows, worth $+$\eur{23,300} in the main scenario; (iv)~adaptive location choice did not beat simple maximin ordering; (v)~New Lock could not reach the quality targets with any number of CPTs.

\textbf{Lessons.} An AI algorithm should be judged by the decision it supports, using an independent audit
and reporting negative results. In investigation planning, site predictability and decision latency
mattered more here than the marginal accuracy of the predictor.
